\section{Introduction}\label{sec:introduction}

The Erd\H{o}s--R\'enyi graph $G(n,p)$ is the random graph on $n$ vertices in
which each edge is present independently with probability $p$.  For a fixed
graph $H$ and fixed $0<p<r<1$, consider the upper-tail event that the
homomorphism density of $H$ in $G(n,p)$ is at least $r^{|E(H)|}$, which is,
to leading order, its typical value in $G(n,r)$.  Two questions
of Chatterjee and Varadhan~\cite{ChatterjeeVaradhanLDP} guide the study of
this event: how small is its probability, and what does $G(n,p)$ typically
look like conditioned on it?  In the dense regime considered here, both
questions are answered by a variational problem on graphons, and the second
question asks for the structure of its minimizers.

These minimizers are known only in special cases.  If the parameters lie in
the replica-symmetric phase, the unique minimizer is the constant graphon, and
the conditioned graph looks like $G(n,r)$.  Lubetzky and
Zhao~\cite{lubetzky2015large} determined this phase exactly when $H$ is
regular.  In the complementary symmetry-breaking phase, the minimizers are
nonconstant, and Lubetzky and Zhao~\cite[Section~7]{lubetzky2015large} asked
whether they are always bipodal, that is, whether the conditioned graph always
splits into two blocks of vertices with homogeneous edge densities within and
between the blocks.  The companion paper~\cite{LimImKimNamYangBipodal}
confirmed this for $d$-regular $H$ near every point of the phase boundary with
target density $r\ne(d-1)/d$, and along a curve approaching the boundary at the
exceptional target density $(d-1)/d$, using the entropy-maximization results of
Kenyon, Radin, Ren, and Sadun~\cite{kenyon2014entropy}.  For nonregular graphs,
the phase boundary is studied in~\cite{LimImKimNamYangBoundary}, where it is
shown, for instance, that stars with many leaves have nonconstant minimizers on
the boundary itself.  Away from the boundary, and for nonregular graphs, even
the following weaker question has remained open: is every minimizer a graphon
with \emph{finitely many} blocks?

A closely related question arises for dense graphs with prescribed edge
density and prescribed density of $H$.  Their typical structure is given by the
graphons of maximal entropy under these two constraints.  Radin, Ren, and
Sadun~\cite{RadinRenSadun2014} and subsequent
works~\cite{KenyonRadinRenSadun2017b,NeemanRadinSadun2023,NeemanRadinSadun2024,RadinSadun2013,RadinSadun2015}
revealed a rich phase diagram when $H$ is a triangle.  The maximizers found
there, numerically and in the regimes where they have been determined
rigorously, are \emph{multipodal}: they have finitely many blocks.  Kenyon,
Radin, Ren, and Sadun~\cite{KenyonRadinRenSadun2017} proved multipodality when
$H$ is a star, and showed in~\cite{kenyon2014entropy} that for every graph $H$
the maximizers are bipodal when the $H$-density is slightly above the
Erd\H{o}s--R\'enyi value, for all but finitely many edge densities.  Beyond
these regimes, multipodality has not been established.

In this paper we show that every minimizer of the upper-tail variational
problem is multipodal, for every graph $H$ and every $0<p<r<1$, and that the
same holds for the maximizers of the entropy with prescribed edge density and
$H$-density, for almost every value of the $H$-density.  In both settings the
blocks are moreover ordered: the blocks can be listed so that the edge density
between two blocks is nondecreasing in the index of each block.

Finite multipodality is not a formal consequence of the fact that only
finitely many densities are prescribed.  Graphons that are uniquely determined
by finitely many subgraph densities can have a very complex structure; see,
for example,~\cite{CooperKralMartins2018,GrzesikKralLovasz2020}.  Our argument
uses both the entropy and the global optimality of the minimizer in an
essential way.

\subsection{Main results}\label{sec:main-results}

A \emph{graphon} is a symmetric measurable function $W:[0,1]^2\to[0,1]$, and
$\W_0$ denotes the space of graphons~\cite{LovaszGraphLimits}.  For a graph $H$
with $m:=|E(H)|\ge1$ edges and a graphon $W$, the edge density, the
homomorphism density of $H$, and the relative entropy with respect to $p$ are
\begin{equation}\label{eq:graphon-quantities}
  e(W)=\int_{[0,1]^2}W,
  \qquad
  t(H,W)=\int_{[0,1]^{V(H)}}\prod_{ij\in E(H)}W(x_i,x_j)\prod_{i\in V(H)}\dd x_i,
  \qquad
  I_p(W)=\int_{[0,1]^2}J_p(W),
\end{equation}
where $J_p(z)=z\log\frac zp+(1-z)\log\frac{1-z}{1-p}$.  For $0<p<r<1$, the
Chatterjee--Varadhan variational problem is
\begin{equation}\label{eq:graphon-variational-problem}
  \Phi_H(p,r)=\inf\bigl\{I_p(W):W\in\W_0,\ t(H,W)\ge r^m\bigr\}.
\end{equation}
Chatterjee and Varadhan~\cite{ChatterjeeVaradhanLDP} proved that the
probability of the event $t(H,G(n,p))\ge r^m$ is
$\exp(-\Phi_H(p,r)\binom n2+o(n^2))$, and that, conditioned on this event,
$G(n,p)$ is close in cut distance to the set of minimizers
of~\eqref{eq:graphon-variational-problem} with probability tending to one.

A graphon $W$ is \emph{multipodal} if there are a finite measurable partition
$[0,1]=A_1\cup\dots\cup A_M$ and numbers $b_{ij}=b_{ji}\in[0,1]$ such that
$W=b_{ij}$ almost everywhere on $A_i\times A_j$.  Its blocks are
\emph{ordered} if they can be labeled so that $b_{ij}\le b_{i'j}$ whenever
$i\le i'$, for every $j$.  Equivalently, the distinct rows $W(x,\cdot)$ of $W$
are totally ordered pointwise.  In graph terms, an ordered multipodal graphon
describes a stochastic block model in which the blocks can be ranked so that
every vertex in a higher-ranked block is at least as likely as every vertex in
a lower-ranked block to be adjacent to any given vertex.  For example, a
bipodal graphon of the form $f(x)f(y)$, where $f$ takes two values, has ordered
blocks; the optimizers constructed in~\cite{LimImKimNamYangBipodal} near the
exceptional target density $(d-1)/d$ are of this form.

\begin{theorem}[Upper tails]\label{thm:upper-tail}
Let $H$ be a graph with at least one edge, and let $0<p<r<1$.  Every minimizer
of~\eqref{eq:graphon-variational-problem} is multipodal.  If $H$ is not a
perfect matching, its blocks are ordered.  If $H$ is a perfect matching, the
unique minimizer is the constant graphon $W\equiv r$.
\end{theorem}

Theorem~\ref{thm:upper-tail} holds throughout the parameter space, including
the symmetry-breaking phase, where the minimizers are not known explicitly.
It does not bound the number of blocks, and it does not assert uniqueness.

We next consider graphons with a prescribed edge density $\eps\in(0,1)$ and a
prescribed $H$-density $\tau$.  Recall the graphon entropy
\[
  s(W)=\int_{[0,1]^2}S_0\bigl(W(x,y)\bigr)\dd x\dd y,
  \qquad
  S_0(z)=-\frac12\bigl(z\log z+(1-z)\log(1-z)\bigr).
\]
Let $\tau_{\max}(\eps):=\max\{t(H,W):e(W)=\eps\}$, and consider
\begin{equation}\label{eq:entropy-maximization}
  \sup\bigl\{s(W):W\in\W_0,\ e(W)=\eps,\ t(H,W)=\tau\bigr\},
  \qquad
  \eps^m<\tau<\tau_{\max}(\eps).
\end{equation}
The maximizers of~\eqref{eq:entropy-maximization} describe the typical
structure of large graphs with edge density $\eps$ and $H$-density
$\tau$~\cite{RadinRenSadun2014,RadinSadun2013}; see
also~\cite{DemboLubetzky2018} for the underlying large-deviation principle.  The two constraints of
this problem can be degenerate at a maximizer, and we need a constraint
qualification (Definition~\ref{def:constraint-qualification}) that allows the
two densities to be corrected independently by bounded perturbations.

\begin{theorem}[Prescribed edge and subgraph densities]\label{thm:fixed-densities}
Let $H$ be a graph with a vertex of degree at least two, let $\eps\in(0,1)$, and
let $\eps^m<\tau<\tau_{\max}(\eps)$.  Every maximizer
of~\eqref{eq:entropy-maximization} that satisfies the constraint qualification
is multipodal with ordered blocks.
\end{theorem}

\begin{cor}\label{cor:almost-every-tau}
Let $H$ and $\eps$ be as in Theorem~\ref{thm:fixed-densities}.  The set of
$\tau\in(\eps^m,\tau_{\max}(\eps))$ for which~\eqref{eq:entropy-maximization} has
a maximizer violating the constraint qualification has Lebesgue measure zero.
Consequently, for almost every such $\tau$, every maximizer is multipodal with
ordered blocks.
\end{cor}

\subsection{Proof strategy}\label{sec:proof-strategy}

The natural starting point is the Euler--Lagrange equation.  For a minimizer
$W$ of~\eqref{eq:graphon-variational-problem}, it states that
$J_p'(W(x,y))=\mu\,\mathcal D_W(x,y)$ for a Lagrange multiplier $\mu>0$, where
$\mathcal D_W$ is the first variation of $t(H,\cdot)$.  One can rewrite this as
the statement that every row $W_x=W(x,\cdot)$ is a critical point of a
functional $\mathcal F_W$ defined on all possible rows.  However, the set of
critical points of a nonlinear functional on an infinite-dimensional space can
be complicated, and the equation alone does not seem to force finitely many
row types for general $H$.  Our proof instead proceeds in three steps.

\begin{enumerate}[label=(\roman*),leftmargin=2.2em]
\item \emph{Every row is a global minimizer.}  Given an arbitrary candidate row
$u:[0,1]\to[0,1]$, we insert into $W$ a new block of vertices of small measure
$\theta$ whose row is $u$, and we then shift the remaining values slightly so
that the constraint holds exactly.  Comparing the cost of this graphon with that
of $W$ to first order in $\theta$ shows that $\mathcal F_W(u)\ge c_W$ for every $u$,
where $c_W$ is an explicit constant.  The functional $\mathcal F_W$ involves the
\emph{rooted densities} of $H$, which record the density of copies of $H$ that
use the new vertex.  Averaging over the rows of $W$ shows that almost every row
attains the bound $c_W$.  Hence all rows of $W$ are global minimizers of one and
the same functional.  See Lemma~\ref{lem:upper-tail-insertion}.

\item \emph{Minimizers are ordered.}  When two rows $u_1,u_2$ are replaced by
their pointwise minimum and maximum, the sum of the entropy and linear parts of
$\mathcal F_W$ over the two rows does not change, while the sum of the rooted
densities strictly increases whenever $u_1$ and $u_2$ cross; this uses a vertex
of $H$ of degree at least two and the fact that $W$ is bounded below by a
positive constant.
Since the multiplier $\mu$ is positive, two crossing minimizers would produce a
pair of rows with strictly smaller total value, which is impossible.  Hence the
set of global minimizers is totally ordered.  See
Proposition~\ref{prop:row-chain}.

\item \emph{There are finitely many minimizers.}  A totally ordered set of
rows could still be a continuum.  The Euler--Lagrange equation for rows is a
compact fixed-point equation, so the set of global minimizers is compact.  The
same strict supermodularity shows that the second variation of $\mathcal F_W$
is coercive on mean-zero perturbations.  This allows us to parametrize the
critical rows near a minimizer analytically by their mean.  If minimizers
accumulated, the Lagrange multiplier for the mean would vanish on an interval,
and a whole interval of means would consist of minimizers; a compactness
argument rules this out.  Finitely many row types then yield finitely many
blocks.  See Theorem~\ref{thm:row-criterion}.
\end{enumerate}

For the problem~\eqref{eq:entropy-maximization} with two constraints, the
same argument applies once both densities can be restored after inserting a
block, which is what the constraint qualification guarantees.  To prove
Corollary~\ref{cor:almost-every-tau}, we show that a maximizer violating the
constraint qualification forces the left difference quotients of the value
function of the problem to be unbounded, which can only happen on a null set
because the value function is monotone.

\subsection{Related work}

The Euler--Lagrange equations of constrained graphon problems have been used
to establish multipodality when they separate.  For stars, the first variation
of the star density is a sum of a function of $x$ and a function of $y$, and
the equation reduces to a scalar equation for the degree function; this is the
mechanism behind the multipodality results of Kenyon, Radin, Ren, and
Sadun~\cite{KenyonRadinRenSadun2017} for edge and star constraints, where a
uniform bound on the number of blocks is also obtained.  Radin and
Sadun~\cite{RadinSadunEdge2Star} studied the optimal graphons of the
edge--$2$-star model in detail.  For the edge--triangle model, Koch~\cite{Koch2016}
studied an ordering of the vertices of optimal graphons, and Neeman, Radin,
and Sadun~\cite{NeemanRadinSadun2023,NeemanRadinSadun2024} established
detailed properties of several phases.  Radin and
Sadun~\cite{RadinSadunNearExtreme} analyzed optimal graphons for constraints
close to the boundary of the feasible region.
For exponential random graph models with nonnegative parameters, which are
Lagrangian relaxations of problems of the above type, Chatterjee and
Diaconis~\cite{ChatterjeeDiaconis2013} showed that the optimizers are constant;
in the symmetry-breaking phase, however, the minimizers
of~\eqref{eq:graphon-variational-problem} are nonconstant and are not
minimizers of the corresponding Lagrangian.

Alvarado, Ramon, and Wang~\cite{AlvaradoRamonWang2023,AlvaradoRamonWang2026}
consider entropy maximization for multi-relational hypergraphons under
constraints on subgraph densities, and state that for almost every value of
the constraints the optimizers are step functions whose number of steps is
bounded in terms of the number of constraints.  Several steps of their
arguments, however, do not hold.  Representatively, the arguments compare step
functions with $k$ and $k+1$ steps and assert that passing to $k+1$ steps
produces no smaller value of the objective.  This fails in general.  In
Appendix~\ref{app:arw} we give constraints on the densities of the edge, the
$2$-star, and the $3$-star, together with a non-convex objective, for which the
bound on the number of steps fails on a set of constraint values of positive
measure; a second example shows the same for the entropy under constraints of
the more general form allowed in~\cite{AlvaradoRamonWang2026}.  The later
preprint also deduces local optimality in $L^1$ from second-order conditions, although
inserting a small block of vertices is a perturbation that is small in $L^1$
and changes the objective at first order.  These examples do not decide whether
the conclusions of~\cite{AlvaradoRamonWang2023,AlvaradoRamonWang2026} hold for
the entropy under subgraph density constraints.  Our proofs are independent of
these works and follow a different route: we never approximate optimizers by
step functions of bounded size, and the insertion of small blocks is precisely
what gives the first-order condition satisfied by every optimizer.
Appendix~\ref{app:arw} contains the counterexamples and explains why they do
not affect the arguments of this paper.

\paragraph{Organization.}
Section~\ref{sec:preliminaries} collects preliminaries, including rooted
densities.  Section~\ref{sec:row-criterion} proves the criterion for finitely
many row types.  Sections~\ref{sec:upper-tail} and~\ref{sec:fixed-densities}
apply it to prove Theorems~\ref{thm:upper-tail} and~\ref{thm:fixed-densities},
and Section~\ref{sec:conclusion} discusses open problems.
Appendix~\ref{app:arw} discusses the preprints~\cite{AlvaradoRamonWang2023,AlvaradoRamonWang2026}.
